% !TeX program = xelatex
% 编译：在 Main 目录运行 latexmk -xelatex main.tex。
% 首页沿用所给 PDF 的 mcmthesis 模板；正文沿用所给 main.tex 的层级。
\documentclass{mcmthesis}

% ===== 基本信息：在此修改年份、队号、题号和论文标题 =====
\newcommand{\ContestYear}{2026}
\newcommand{\TeamNumber}{000000} % 待替换为本次参赛队号
\newcommand{\ProblemLetter}{B}
\newcommand{\PaperTitle}{Sand Control in Taklamakan Desert} % 待替换为论文标题
\mcmsetup{
  tstyle=\color{red}\bfseries,
  tcn=\TeamNumber, problem=\ProblemLetter,
  sheet=true, titleinsheet=true, keywordsinsheet=true,
  titlepage=false, abstract=true
}
\renewcommand{\headset}{\ContestYear\\MCM/ICM\\Summary Sheet}
\title{\PaperTitle}

\usepackage{times}
\usepackage{indentfirst}
\usepackage{needspace,tikz,subcaption,float,placeins}
\usetikzlibrary{arrows.meta,positioning}
\usepackage{lastpage}
% 中英文文献库：XeLaTeX + BibTeX；正文使用 \cite{REF01} 等引用键。
\usepackage{xeCJK}
\setmainfont{texgyretermes-regular.otf}[BoldFont=texgyretermes-bold.otf,ItalicFont=texgyretermes-italic.otf,BoldItalicFont=texgyretermes-bolditalic.otf]
\setCJKmainfont{FandolSong-Regular.otf}[BoldFont=FandolSong-Bold.otf,ItalicFont=FandolKai-Regular.otf]
\setCJKsansfont{FandolHei-Regular.otf}[BoldFont=FandolHei-Bold.otf]
\setCJKmonofont{FandolFang-Regular.otf}
\usepackage[backend=bibtex,style=numeric,sorting=none,maxbibnames=3,giveninits=true]{biblatex}
\usepackage{xurl}
\addbibresource{bibliography/references.bib}
% 有 DOI 时不重复输出同一文献的网址。
\AtEveryBibitem{\iffieldundef{doi}{}{\clearfield{url}}}
\graphicspath{{Figure/}}
\geometry{headheight=14pt}
\setlength{\emergencystretch}{1.2em}
\hypersetup{hidelinks,pdfauthor={},pdftitle={MCM/ICM Paper}}

% 保留旧论文的符号表、公式分组和交替灰底表格工具。
\newcommand{\shade}{\rowcolor{gray!22}}
\newcommand{\modelvars}[1]{%
  \par\begingroup\small
  \setlength{\tabcolsep}{4pt}\renewcommand{\arraystretch}{1.10}
  \noindent\begin{tabularx}{\linewidth}{@{}>{\raggedright\arraybackslash}p{0.29\linewidth} >{\raggedright\arraybackslash}X@{}}
  \toprule\textbf{Symbol} & \textbf{Meaning} \\\midrule
  #1
  \bottomrule\end{tabularx}\par\endgroup
}
\newenvironment{modeleq}
  {\par\addvspace{0.4\baselineskip}\noindent\begin{minipage}{\linewidth}\setlength{\abovedisplayskip}{6pt}\setlength{\belowdisplayskip}{6pt}}
  {\end{minipage}\par\addvspace{0.4\baselineskip}}

% 占位内容仅供检查排版，写作时逐一替换；不包含旧题的模型与结论。
\newcommand{\placeholder}[1]{\par\noindent\textit{[#1]}\par}

\begin{document}

% ===== 第一页：Summary Sheet（摘要不编入正文章节编号） =====
% mcmthesis 要求先声明 abstract 和 keywords，再调用 maketitle。
\begin{abstract}
\placeholder{Write the problem overview, modeling methods, principal numerical
results, sensitivity findings, and recommendations here. Keep the complete
summary and keywords on this first page.}
\begin{keywords}
Keyword 1; Keyword 2; Keyword 3; Keyword 4
\end{keywords}
\end{abstract}
\hypersetup{pageanchor=false}
\maketitle
\hypersetup{pageanchor=true}

% ===== 第二页起：目录及正文 =====
\pagestyle{main}
\tableofcontents
\clearpage

\section{Introduction}
\subsection{Problem background}
The Taklamankan desert covers more then 337,600 square kilometers, which is the largest desert in China. And in this rural area, Chinese government have set up one of the biggest "Green Sand Protective Belt", with total length about 3,046 kilometers.
However, sand control is not a static policy, but a \textbf{cohesive and dynamic} system with a list of parameter, including the distribution of water, funding, labor and temporary diaster.
This requires us to develop a well-rounded and long-term method to quantify the condition of sand protective belt.


The main factors affecting the protective belt are summarized in
\textbf{Table~\ref{tab:protection-factors}.}
\begin{table}[htbp]
  \centering
  \caption{Key factors and quantities in protective-belt management.}
  \label{tab:protection-factors}
  \begingroup\small
  \renewcommand{\arraystretch}{1.18}
  \begin{tabularx}{\linewidth}{@{}>{\raggedright\arraybackslash}p{0.17\linewidth}
    >{\raggedright\arraybackslash}p{0.37\linewidth}
    >{\raggedright\arraybackslash}X@{}}
    \toprule
    \textbf{Factor} & \textbf{Representative quantities} & \textbf{Role in belt protection} \\
    \midrule
    Water supply & Ecological replenishment volume (m\textsuperscript{3}/year); irrigation timing
      & Sustains soil moisture and vegetation growth under water constraints. \\
    Funding & Total budget (CNY); investment shares in construction, O\&M, and supporting facilities
      & Determines the scale of new treatment and the resources available for continued operation. \\
    Labor & Annual labor input (person-days/year); seasonal allocation
      & Supports construction, irrigation, maintenance, and post-disaster repair. \\
    Disaster shocks & Event frequency (events/year), duration (days), and severity
      & Causes vegetation loss and changes soil moisture and protective performance. \\
    \bottomrule
  \end{tabularx}
  \endgroup
\end{table}

\subsection{Problem restatement}
We aim to design a protective belt that remains ecologically effective and
financially sustainable. For a representative section along the Taklamakan Desert
margin, our objectives are to:
\begin{enumerate}
  \item \textbf{Establish a coupled dynamic model.}
  Use relevant data to describe \textbf{vegetation coverage},
  \textbf{effective soil moisture}, and their \textbf{interactions with wind and sand}.

  \item \textbf{Optimize construction and resource allocation.}
  Balance ecological benefits against economic costs under a given budget and
  planning horizon (recommended: 10--20 years). Optimize the
  \textbf{annual newly treated area}, \textbf{water replenishment volume and timing},
  \textbf{investment shares} for construction, O\&M, and revenue-generating facilities
  such as photovoltaics, and \textbf{labor input}. Test \textbf{investment sensitivity}
  and propose an \textbf{actionable process criterion}.

  \item \textbf{Assess resilience and long-term benefits.}
  Incorporate \textbf{stochastic disasters} (sandstorms, droughts, wind erosion,
  and sand burial) and \textbf{climate-driven changes} in wind and moisture conditions.
  Evaluate \textbf{resistance and recovery} and project
  \textbf{long-term ecological and economic benefits}.

  \item \textbf{Provide practical recommendations.}
  Summarize feasibility, implementation priorities, fiscal burdens, and expected
  benefits in a one-page brief for authorities and the public.
\end{enumerate}
\subsection{Our work}
\placeholder{Describe the overall framework and the links among the models.}
% \begin{figure}[H]\centering
% \includegraphics[width=\linewidth]{workflow.pdf}
% \caption{Overall modeling framework.}\label{fig:framework}
% \end{figure}

\clearpage
\section{Assumptions and Justifications}
\begin{enumerate}
  \item \textbf{Assumption:} \textit{[State an assumption.]}

  \textbf{Justification:} \textit{[Explain why it is reasonable.]}
  % 按上述格式继续添加 \item，每条假设后给出对应解释。
\end{enumerate}

\clearpage
\section{Notations}
\placeholder{State the units and numerical conventions used throughout the paper.}
\subsection*{Sets and indices}
\modelvars{\textit{[Index]} & \textit{[Definition]} \\}
\subsection*{Variables and derived quantities}
\modelvars{\textit{[Variable]} & \textit{[Meaning and unit]} \\}
\subsection*{Baseline quantities}
\modelvars{\textit{[Parameter]} & \textit{[Value, unit, and source]} \\}

\clearpage
\section{Model 1: Model formulation}\label{sec:model1}
\subsection{Data and model preparation}
\placeholder{Introduce data sources, preprocessing, units, variables, and parameters.}
\subsection{Model construction}
\placeholder{Develop the governing equations, resource balances, feedbacks,
performance indicators, objective function, and constraints.}
\subsection{Solution and results}
\placeholder{Describe the solution procedure, interpret the recommended strategy
and resource allocation, and evaluate long-term performance.}

% 第 4、5、6 节分别为 Model 1、Model 2、Model 3。
\clearpage
\section{Model 2: Resource allocation and system feedback}\label{sec:model2}
\subsection{From resources to system improvements}
\placeholder{Develop the allocation model and report its results.}
\subsection{How investment feeds back into the system}
\placeholder{Quantify the effects of allocation on the core model.}

\clearpage
\section{Model 3: Sensitivity and preference robustness}\label{sec:model3}
\subsection{Sensitivity analysis}
\placeholder{Perturb uncertain parameters individually and jointly, and explain
the mechanisms behind the sensitivity results.}
\subsection{Robustness analysis}
\placeholder{Compare alternative preferences and robust strategies.}

\clearpage
\phantomsection
\addcontentsline{toc}{section}{References}
\begin{thebibliography}{99}\small
\bibitem{source-placeholder}
\textit{[Replace this entry with an actual cited source.]}
\end{thebibliography}

\end{document}
